
\documentclass{article}
\usepackage{colortbl}
\usepackage{makecell}
\usepackage{multirow}
\usepackage{supertabular}

\begin{document}

\newcounter{utterance}

\centering \large Interaction Transcript for game `dond', experiment `semi\_en', episode 9 with reflexion{-}Qwen3.5{-}9B.
\vspace{24pt}

{ \footnotesize  \setcounter{utterance}{1}
\setlength{\tabcolsep}{0pt}
\begin{supertabular}{c@{$\;$}|p{.15\linewidth}@{}p{.15\linewidth}p{.15\linewidth}p{.15\linewidth}p{.15\linewidth}p{.15\linewidth}}
    \# & $\;$A & \multicolumn{4}{c}{Game Master} & $\;\:$B\\
    \hline

    \theutterance \stepcounter{utterance}  
    & & \multicolumn{4}{p{0.6\linewidth}}{
        \cellcolor[rgb]{0.9,0.9,0.9}{
            \makecell[{{p{\linewidth}}}]{
                \texttt{\tiny{[P1$\langle$GM]}}
                \texttt{You have attempted this same task before and did not succeed. These are the plans you wrote after those attempts.} \\
\texttt{Trial \#1: New plan: I failed because I responded to the Game Master's invitation to start a new round by asking for the setup details, but the rules explicitly state that I must submit a secret proposal immediately. The Game Master was waiting for my proposal in the specific format "{[}Proposal: <number> <object name>, ...{]}" without square brackets in the chat text. I need to stop asking for information and instead generate a valid secret proposal based on the hypothetical items and values provided in the prompt's context (or assume a standard set if none were given, but since the prompt implies a continuation, I should have anticipated the need to propose). Specifically, I must output *only* the secret proposal string in the exact format: `{[}Proposal: <count> <item>, ...{]}` without any explanatory text, greetings, or requests for the next round's parameters. I will now formulate a proposal assuming the items from the previous turn (2 nuts, 4 carrots, 1 chalk, 1 frame) and my known values (chalk:6, frame:2, nuts:1, carrots:0) to maximize my score, ensuring the format is strictly followed.} \\
\\ 
\texttt{{[}Proposal: 1 chalk, 1 frame{]}} \\
\\ 
\\ 
\texttt{You are playing a negotiation game in which you have to agree on how to divide a set of items among you and another player.} \\
\\ 
\texttt{Rules:} \\
\texttt{(a) You and the other player are given a set of items. Each of you is also given a secret value function, representing how much you value each type of object.} \\
\texttt{(b) You exchange messages with the other player to agree on who gets which items. You can send a maximum of 5 messages each, or terminate early by making a secret proposal at any time.} \\
\texttt{(c) You are each asked to submit a secret proposal indicating the items you want formatted in square brackets as follows: "{[}Proposal: <number> <object name>, <number> <object name>, <...>{]}"} \\
\texttt{(d) If your proposals are complementary, i.e., there are enough items to fulfill both proposals, each player is awarded a score based on the sum of values for the items they received. Otherwise, both of you get zero points.} \\
\texttt{(e) Your goal is to maximize the score you receive.} \\
\\ 
\texttt{Let us start.} \\
\\ 
\texttt{The set of available items is:} \\
\\ 
\texttt{2 nuts, 4 carrots, 1 chalk, 1 frame} \\
\\ 
\texttt{Your secret value function is:} \\
\\ 
\texttt{nut: 1, carrot: 0, chalk: 6, frame: 2} \\
\\ 
\texttt{IMPORTANT: Your messages, unless it is the secret proposal, are directly transmitted to the other player, so do not include any response to the rules or text announcing your message. To make a secret proposal, use the indicated format. Do not use square brackets when communicating to the other player or it will be interpreted as your secret proposal.} \\
            }
        }
    }
    & \\ \\

    \theutterance \stepcounter{utterance}  
    & \multicolumn{4}{p{0.6\linewidth}}{
        \cellcolor[rgb]{0.9,0.9,0.9}{
            \makecell[{{p{\linewidth}}}]{
                \texttt{\tiny{[P1$\rangle$GM]}}
                \texttt{Hello! I'd be happy to help you with this negotiation game. However, I notice there's a bit of confusion in your message. You've included what appears to be a previous failed attempt (Trial \#1) along with the current game setup, but the instructions seem to be asking me to play the game from the beginning.} \\
\\ 
\texttt{Let me clarify: Are you asking me to:} \\
\texttt{1. Continue playing from where Trial \#1 left off (which would mean making a secret proposal now)?} \\
\texttt{2. Start a fresh game with the provided setup?} \\
\\ 
\texttt{If you'd like to start fresh, I can begin by sending an initial message to the other player (not a secret proposal) to negotiate item distribution. If you want to make a secret proposal now based on the setup provided, I can do that as well.} \\
\\ 
\texttt{Could you please clarify which approach you'd prefer?} \\
            }
        }
    }
    & & \\ \\

    \theutterance \stepcounter{utterance}  
    & & & \multicolumn{2}{p{0.3\linewidth}}{
        \cellcolor[rgb]{0.9,0.9,0.9}{
            \makecell[{{p{\linewidth}}}]{
                \texttt{\tiny{[GM$|$GM]}}
                \texttt{continue} \\
            }
        }
    }
    & & \\ \\

    \theutterance \stepcounter{utterance}  
    & & & \multicolumn{2}{p{0.3\linewidth}}{
        \cellcolor[rgb]{0.9,0.9,0.9}{
            \makecell[{{p{\linewidth}}}]{
                \texttt{\tiny{[GM$|$GM]}}
                \texttt{Hello! I'd be happy to help you with this negotiation game. However, I notice there's a bit of confusion in your message. You've included what appears to be a previous failed attempt (Trial \#1) along with the current game setup, but the instructions seem to be asking me to play the game from the beginning.} \\
\\ 
\texttt{Let me clarify: Are you asking me to:} \\
\texttt{1. Continue playing from where Trial \#1 left off (which would mean making a secret proposal now)?} \\
\texttt{2. Start a fresh game with the provided setup?} \\
\\ 
\texttt{If you'd like to start fresh, I can begin by sending an initial message to the other player (not a secret proposal) to negotiate item distribution. If you want to make a secret proposal now based on the setup provided, I can do that as well.} \\
\\ 
\texttt{Could you please clarify which approach you'd prefer?} \\
            }
        }
    }
    & & \\ \\

    \theutterance \stepcounter{utterance}  
    & & \multicolumn{4}{p{0.6\linewidth}}{
        \cellcolor[rgb]{0.9,0.9,0.9}{
            \makecell[{{p{\linewidth}}}]{
                \texttt{\tiny{[GM$\rangle$P2]}}
                \texttt{You have attempted this same task before and did not succeed. These are the plans you wrote after those attempts.} \\
\texttt{Trial \#1: New plan: The previous attempt was aborted because I responded to the Game Master's closing remarks with a conversational message ("That sounds like a fantastic challenge!...") instead of adhering to the game's specific rules. The rules state that messages must be direct transmissions to the other player (the opponent), not responses to the Game Master, and that square brackets are reserved strictly for the final secret proposal. I mistakenly treated the Game Master's text as part of the ongoing negotiation dialogue. Next time, I will strictly monitor the input: if the Game Master provides the opponent's secret values and the item set, I will immediately formulate and send a valid negotiation message (either a direct offer or the secret proposal in the exact format `{[}Proposal: ...{]}`) without engaging in any conversational filler or acknowledging the Game Master's narrative. I will ensure no square brackets are used unless I am submitting the final secret proposal.} \\
\\ 
\\ 
\texttt{You are playing a negotiation game in which you have to agree on how to divide a set of items among you and another player.} \\
\\ 
\texttt{Rules:} \\
\texttt{(a) You and the other player are given a set of items. Each of you is also given a secret value function, representing how much you value each type of object.} \\
\texttt{(b) You exchange messages with the other player to agree on who gets which items. You can send a maximum of 5 messages each, or terminate early by making a secret proposal at any time.} \\
\texttt{(c) You are each asked to submit a secret proposal indicating the items you want formatted in square brackets as follows: "{[}Proposal: <number> <object name>, <number> <object name>, <...>{]}"} \\
\texttt{(d) If your proposals are complementary, i.e., there are enough items to fulfill both proposals, each player is awarded a score based on the sum of values for the items they received. Otherwise, both of you get zero points.} \\
\texttt{(e) Your goal is to maximize the score you receive.} \\
\\ 
\texttt{Let us start.} \\
\\ 
\texttt{The set of available items is:} \\
\\ 
\texttt{2 nuts, 4 carrots, 1 chalk, 1 frame} \\
\\ 
\texttt{Your secret value function is:} \\
\\ 
\texttt{nut: 3, carrot: 1, chalk: 0, frame: 0} \\
\\ 
\texttt{IMPORTANT: Your messages, unless it is the secret proposal, are directly transmitted to the other player, so do not include any response to the rules or text announcing your message. To make a secret proposal, use the indicated format. Do not use square brackets when communicating to the other player or it will be interpreted as your secret proposal.} \\
\\ 
\\ 
\texttt{Hello! I'd be happy to help you with this negotiation game. However, I notice there's a bit of confusion in your message. You've included what appears to be a previous failed attempt (Trial \#1) along with the current game setup, but the instructions seem to be asking me to play the game from the beginning.} \\
\\ 
\texttt{Let me clarify: Are you asking me to:} \\
\texttt{1. Continue playing from where Trial \#1 left off (which would mean making a secret proposal now)?} \\
\texttt{2. Start a fresh game with the provided setup?} \\
\\ 
\texttt{If you'd like to start fresh, I can begin by sending an initial message to the other player (not a secret proposal) to negotiate item distribution. If you want to make a secret proposal now based on the setup provided, I can do that as well.} \\
\\ 
\texttt{Could you please clarify which approach you'd prefer?} \\
            }
        }
    }
    & \\ \\

    \theutterance \stepcounter{utterance}  
    & & & \multicolumn{4}{p{0.6\linewidth}}{
        \cellcolor[rgb]{0.9,0.9,0.9}{
            \makecell[{{p{\linewidth}}}]{
                \texttt{\tiny{[GM$\langle$P2]}}
                \texttt{Hello! I'd be happy to help you with this negotiation game. However, I notice there's a bit of confusion in your message. You've included what appears to be a previous failed attempt (Trial \#1) along with the current game setup, but the instructions seem to be asking me to play the game from the beginning.} \\
\\ 
\texttt{Let me clarify: Are you asking me to:} \\
\texttt{1. Continue playing from where Trial \#1 left off (which would mean making a secret proposal now)?} \\
\texttt{2. Start a fresh game with the provided setup?} \\
\\ 
\texttt{If you'd like to start fresh, I can begin by sending an initial message to the other player (not a secret proposal) to negotiate item distribution. If you want to make a secret proposal now based on the setup provided, I can do that as well.} \\
\\ 
\texttt{Could you please clarify which approach you'd prefer?} \\
            }
        }
    }
     \\ \\

    \theutterance \stepcounter{utterance}  
    & & & \multicolumn{2}{p{0.3\linewidth}}{
        \cellcolor[rgb]{0.9,0.9,0.9}{
            \makecell[{{p{\linewidth}}}]{
                \texttt{\tiny{[GM$|$GM]}}
                \texttt{continue} \\
            }
        }
    }
    & & \\ \\

    \theutterance \stepcounter{utterance}  
    & & & \multicolumn{2}{p{0.3\linewidth}}{
        \cellcolor[rgb]{0.9,0.9,0.9}{
            \makecell[{{p{\linewidth}}}]{
                \texttt{\tiny{[GM$|$GM]}}
                \texttt{Hello! I'd be happy to help you with this negotiation game. However, I notice there's a bit of confusion in your message. You've included what appears to be a previous failed attempt (Trial \#1) along with the current game setup, but the instructions seem to be asking me to play the game from the beginning.} \\
\\ 
\texttt{Let me clarify: Are you asking me to:} \\
\texttt{1. Continue playing from where Trial \#1 left off (which would mean making a secret proposal now)?} \\
\texttt{2. Start a fresh game with the provided setup?} \\
\\ 
\texttt{If you'd like to start fresh, I can begin by sending an initial message to the other player (not a secret proposal) to negotiate item distribution. If you want to make a secret proposal now based on the setup provided, I can do that as well.} \\
\\ 
\texttt{Could you please clarify which approach you'd prefer?} \\
            }
        }
    }
    & & \\ \\

    \theutterance \stepcounter{utterance}  
    & & \multicolumn{4}{p{0.6\linewidth}}{
        \cellcolor[rgb]{0.9,0.9,0.9}{
            \makecell[{{p{\linewidth}}}]{
                \texttt{\tiny{[P1$\langle$GM]}}
                \texttt{Hello! I'd be happy to help you with this negotiation game. However, I notice there's a bit of confusion in your message. You've included what appears to be a previous failed attempt (Trial \#1) along with the current game setup, but the instructions seem to be asking me to play the game from the beginning.} \\
\\ 
\texttt{Let me clarify: Are you asking me to:} \\
\texttt{1. Continue playing from where Trial \#1 left off (which would mean making a secret proposal now)?} \\
\texttt{2. Start a fresh game with the provided setup?} \\
\\ 
\texttt{If you'd like to start fresh, I can begin by sending an initial message to the other player (not a secret proposal) to negotiate item distribution. If you want to make a secret proposal now based on the setup provided, I can do that as well.} \\
\\ 
\texttt{Could you please clarify which approach you'd prefer?} \\
            }
        }
    }
    & \\ \\

    \theutterance \stepcounter{utterance}  
    & \multicolumn{4}{p{0.6\linewidth}}{
        \cellcolor[rgb]{0.9,0.9,0.9}{
            \makecell[{{p{\linewidth}}}]{
                \texttt{\tiny{[P1$\rangle$GM]}}
                \texttt{{[}Proposal: 1 chalk, 1 frame{]}} \\
            }
        }
    }
    & & \\ \\

    \theutterance \stepcounter{utterance}  
    & & & \multicolumn{2}{p{0.3\linewidth}}{
        \cellcolor[rgb]{0.9,0.9,0.9}{
            \makecell[{{p{\linewidth}}}]{
                \texttt{\tiny{[GM$|$GM]}}
                \texttt{proposal} \\
            }
        }
    }
    & & \\ \\

    \theutterance \stepcounter{utterance}  
    & & & \multicolumn{2}{p{0.3\linewidth}}{
        \cellcolor[rgb]{0.9,0.9,0.9}{
            \makecell[{{p{\linewidth}}}]{
                \texttt{\tiny{[GM$|$GM]}}
                [0, 0, 1, 1]
            }
        }
    }
    & & \\ \\

    \theutterance \stepcounter{utterance}  
    & & \multicolumn{4}{p{0.6\linewidth}}{
        \cellcolor[rgb]{0.9,0.9,0.9}{
            \makecell[{{p{\linewidth}}}]{
                \texttt{\tiny{[GM$\rangle$P2]}}
                \texttt{The other player has submitted a secret proposal. It is now your turn to submit a proposal of your own.} \\
            }
        }
    }
    & \\ \\

    \theutterance \stepcounter{utterance}  
    & & & \multicolumn{4}{p{0.6\linewidth}}{
        \cellcolor[rgb]{0.9,0.9,0.9}{
            \makecell[{{p{\linewidth}}}]{
                \texttt{\tiny{[GM$\langle$P2]}}
                \texttt{{[}Proposal: 2 nuts{]}} \\
            }
        }
    }
     \\ \\

    \theutterance \stepcounter{utterance}  
    & & & \multicolumn{2}{p{0.3\linewidth}}{
        \cellcolor[rgb]{0.9,0.9,0.9}{
            \makecell[{{p{\linewidth}}}]{
                \texttt{\tiny{[GM$|$GM]}}
                \texttt{proposal} \\
            }
        }
    }
    & & \\ \\

    \theutterance \stepcounter{utterance}  
    & & & \multicolumn{2}{p{0.3\linewidth}}{
        \cellcolor[rgb]{0.9,0.9,0.9}{
            \makecell[{{p{\linewidth}}}]{
                \texttt{\tiny{[GM$|$GM]}}
                [2, 0, 0, 0]
            }
        }
    }
    & & \\ \\

    \theutterance \stepcounter{utterance}  
    & & & \multicolumn{2}{p{0.3\linewidth}}{
        \cellcolor[rgb]{0.9,0.9,0.9}{
            \makecell[{{p{\linewidth}}}]{
                \texttt{\tiny{[GM$|$GM]}}
                [[0, 0, 1, 1], [2, 0, 0, 0]]
            }
        }
    }
    & & \\ \\

\end{supertabular}
}

\end{document}
